% Section content template for individual Educative lessons
% This file contains the actual content for a specific section
% Generated from Educative API section content

% Section: Code for the Vending Machine
% Section ID: 5842069681864704
% Section Slug: code-for-the-vending-machine
% Generated: 2025-12-11T20:26:25.297332

We've reviewed different aspects of the vending machine problem and observed the attributes attached to the problem using various UML diagrams. Let us now explore the more practical side of things, where we will work on implementing the vending machine using multiple languages. This is usually the last step in an object-oriented design interview process.
\\
We have chosen the following languages to write the skeleton code of the different classes present in the vending machine:

\begin{itemize}
\item
\protect\phantomsection\label{xxqLi_iRC1y_E4voVN1LZ}
Java
\item
\protect\phantomsection\label{2prKR2E0HRs-7MH2KES1T}
C\#
\item
\protect\phantomsection\label{iovjOaMCCGVy7usCBmt29}
Python
\item
\protect\phantomsection\label{mts-Gd4wrtYXKOna3C9un}
C++
\item
\protect\phantomsection\label{xreHb10shtd1z_K9PGo0j}
JavaScript
\end{itemize}
\\
If you want to skip directly to the implementation and see the complete workflow, simply scroll down or click \hyperref[Executable-code-Vending-machine]{Executable Code: Vending Machine system}.

\subsection{Vending machine classes}\label{QHN0J5Js0HfV6qMIBi3BG}
\\
This section will provide the skeleton code of the classes designed in the class diagram lesson.
\begin{quote}
\textbf{Note:} For simplicity, we are not defining getter and setter functions. The reader can assume that all class attributes are private, accessed through their respective public getter methods, and modified only through their public method functions.
\end{quote}
\subsubsection{Enumerations}\label{VyVrBLUQl6c8XvpAPCVIt}
\\
The following code defines the enumeration used in the vending machine system:
\begin{quote}
\textbf{Note:} JavaScript does not support enumerations, so we will use the\ Object.freeze() method as an alternative that freezes an object and prevents further modifications.
\end{quote}

\textbf{Enum definitions}

\begin{lstlisting}[language=Java]
public enum ProductType {
    CHOCOLATE, SNACK, BEVERAGE, OTHER
}
\end{lstlisting}

\subsubsection{State}\label{d6DcwCJ0jKUAiX5iUpOvw-2}
\\
The \texttt{State} interface declares the two user-facing operations whose behavior varies by the machine's current mode. Each concrete state class implements these methods, but leaves their bodies empty here as placeholders for the real logic in \texttt{VendingMachine}.

\textbf{The State interface and its derived classes}

\begin{lstlisting}[language=Java]
// State.java
public interface State {
    /** Called when the user inserts money. */
    void insertMoney(VendingMachine vm, double amount);

    /** Called when the user selects a rack. */
    void selectProduct(VendingMachine vm, int rackNumber);
}

// NoMoneyInsertedState.java
public class NoMoneyInsertedState implements State {
    @Override
    public void insertMoney(VendingMachine vm, double amount) {
        // to be implemented: record amount and switch to MoneyInsertedState
    }

    @Override
    public void selectProduct(VendingMachine vm, int rackNumber) {
        // to be implemented: reject selection when no money
    }
}

// MoneyInsertedState.java
public class MoneyInsertedState implements State {
    @Override
    public void insertMoney(VendingMachine vm, double amount) {
        // to be implemented: add to current balance
    }

    @Override
    public void selectProduct(VendingMachine vm, int rackNumber) {
        // to be implemented: check availability and payment, then dispense or refund
    }
}

// DispenseState.java
public class DispenseState implements State {
    @Override
    public void insertMoney(VendingMachine vm, double amount) {
        // to be implemented: defer until dispensing completes
    }

    @Override
    public void selectProduct(VendingMachine vm, int rackNumber) {
        // to be implemented: defer until dispensing completes
    }
}
\end{lstlisting}

\subsubsection{Product, rack, and inventory}\label{d6DcwCJ0jKUAiX5iUpOvw-3}
\\
At the heart of our vending-machine model are three classes that work together to represent what's in each slot and where:

\begin{itemize}
\item
\protect\phantomsection\label{Ol2vlvq_TWlACV2RDtHUN}
\texttt{Product} encapsulates the details of an item for sale---its unique ID, display name, price, and category (chocolate, snack, beverage, or other).
\item
\protect\phantomsection\label{jM0p6t-AOB4zqlYAzEupx}
\texttt{Rack} represents a single slot in the machine. It knows which \texttt{Product} it holds and how many units remain. When you load or dispense stock, you do so by interacting with the appropriate rack.
\item
\protect\phantomsection\label{hIkjLHjrMKFlY4I7IBgso}
\texttt{Inventory} serves as the lookup service for every rack. It maintains a map from rack numbers to \texttt{Rack} instances, so when the machine needs to check availability or list its contents, it simply asks the \texttt{Inventory} for rack \# N (or for a list of all racks).
\end{itemize}
\\
Together, these classes cleanly separate ``what the items are'' (\texttt{Product}), ``where they live'' (\texttt{Rack}), and ``how to find them'' (\texttt{Inventory}).
\\
The definition of these classes is provided below:

\textbf{The Product, Rack, and Inventory classes}

\begin{lstlisting}[language=Java]
// ===== Product =====
public class Product {
    private final String name;
    private final int id;
    private final double price;
    private final ProductType type;

    public Product(int id, String name, double price, ProductType type) {
        // constructor implementation here
    }

    public int getId() { return 0; }
    public String getName() { return null; }
    public double getPrice() { return 0.0; }
    public ProductType getType() { return null; }
}

// ===== Rack =====
public class Rack {
    private final int rackNumber;
    private Product product;
    private int quantity;

    public Rack(int rackNumber) {
        // constructor implementation here
    }

    public int getRackNumber() { return 0; }
    public boolean isEmpty() { return false; }
    public void loadProduct(Product product, int qty) { }
    public Product peekProduct() { return null; }
    public void dispenseOne() { }
}

// ===== Inventory =====
import java.util.Collection;

public class Inventory {
    public void addRack(Rack rack) { }
    public Rack getRack(int rackNumber) { return null; }
    public Collection<Rack> allRacks() { return null; }
}
\end{lstlisting}

\subsubsection{Vending machine}\label{d6DcwCJ0jKUAiX5iUpOvw-4}
\\
The \texttt{VendingMachine} class is the central controller of our system. Implemented as a Singleton, it encapsulates:

\begin{itemize}
\item
\protect\phantomsection\label{I2_LmUI-fcO3WTcigp0UY}
It creates three State instances (no money, money inserted, dispense) up front, between which it switches.
\item
\protect\phantomsection\label{q9I7CCw1bQ6Q71KOx5y3m}
\texttt{currentState}: The active state object.
\item
\protect\phantomsection\label{PK4RzOQAx882gqYxVewah}
\texttt{currentAmount}: Running total of cash the user has put in.
\item
\protect\phantomsection\label{XiMzv8woBmEbOI9Un4OF0}
\texttt{selectedRack}: The rack number the user has picked (--1 if none).
\item
\protect\phantomsection\label{bO9j4UbkMOlC2W1TaqXfx}
\texttt{inventory}: The Inventory instance that owns all rack objects.
\end{itemize}
\\
All user calls (\texttt{insertMoney} and \texttt{selectProduct}) are simply delegated to the \texttt{currentState}. The machine handles the actual dispensing/refunding and resets to the ``No money'' state when each transaction completes. Administrative methods (\texttt{addRack}/\texttt{loadProduct}) sit on the context here, not in the State interface.
\\
The definition of this class is given below:

\textbf{The VendingMachine class}

\begin{lstlisting}[language=Java]
public class VendingMachine {
    private final State noMoneyState = new NoMoneyInsertedState();
    private final State moneyInsertedState = new MoneyInsertedState();
    private final State dispenseState = new DispenseState();

    private State currentState = noMoneyState;
    private double currentAmount = 0.0;
    private int selectedRack = -1;
    private final Inventory inventory = new Inventory();

    private static VendingMachine instance;
    private VendingMachine() {}

    public static VendingMachine getInstance() {
        if (instance == null) {
            instance = new VendingMachine();
        }
        return instance;
    }

    State getNoMoneyState() { return noMoneyState; }
    State getMoneyInsertedState() { return moneyInsertedState; }
    State getDispenseState() { return dispenseState; }

    void setState(State state) { this.currentState = state; }

    double getCurrentAmount() { return currentAmount; }
    void addToCurrentAmount(double amt) { /* to be implemented */ }
    Inventory getInventory() { return inventory; }
    void setSelectedRack(int rack) { /* to be implemented */ }

    public void insertMoney(double amount) {
        currentState.insertMoney(this, amount);
    }

    public void selectProduct(int rackNumber) {
        currentState.selectProduct(this, rackNumber);
    }

    public void dispenseProduct() { /* to be implemented */ }
    public void refund() { /* to be implemented */ }
    public void addRack(Rack rack) { /* to be implemented */ }
    public void loadProduct(int rackNumber, Product product, int qty) { /* to be implemented */ }
    public void showInventory() { /* to be implemented */ }
}
\end{lstlisting}

\subsection{Executable code: Vending machine}\label{h5Cmrjo5dJvEFleiCRqna}
\\
Below is a fully self-contained, runnable program in Java, C\#, C++, Python, and JavaScript that demonstrates the core workflows of the Vending Machine system. The main driver code of the system resides in the~\texttt{Driver.java},~\texttt{Driver.cs},~\texttt{Driver.py},~\texttt{Driver.cpp}, and~\texttt{Driver.js}~for all respective languages. You can click the ``Run'' button to execute the code.

\subsubsection{What does this code show}\label{l0n-xNKms8P4u_4lm29am}

\begin{itemize}
\item
\protect\phantomsection\label{Xjotngic98NVJyf3lRQMJ}
\textbf{System initialization:} Sets up the vending machine with racks, products, and their quantities.
\item
\protect\phantomsection\label{OTMypr7kMMUtC4sDeZ3No}
\textbf{Scenario 1 (Customer buys a product):} The customer inserts money, selects a product using the rack number, and either receives the product and change (if any) or an error message (if payment is insufficient).
\item
\protect\phantomsection\label{JbkvtawUh4jn9koZC-jR8}
\textbf{Scenario 2 (Operator adds/removes a product):} The operator restocks or removes a product from the machine.
\item
\protect\phantomsection\label{ZGyx8xzhVSmvX-WIyXId4}
\textbf{Scenario 3 (Handling edge cases):} Demonstrates behaviors such as underpayment, exact payment, and overpayment, as well as selecting an empty rack.
\end{itemize}
\\
\textbf{Hint: Try customizing the code!}
\\
This code is designed for experimentation, not just reading. You can edit, extend, or modify any part of the example.

\textbf{Executable Code: Vending Machine}

\begin{lstlisting}[language=Java]
public class Driver {
    public static void main(String[] args) {
        VendingMachine vm = VendingMachine.getInstance();

        // Setup
        vm.addRack(new Rack(1)); vm.addRack(new Rack(2)); vm.addRack(new Rack(3));
        Product choc = new Product(101, "Chocolate Bar", 1.50, ProductType.CHOCOLATE);
        Product snack= new Product(102, "Potato Chips", 2.00, ProductType.SNACK);
        Product bev  = new Product(103, "Soda Can",      2.50, ProductType.BEVERAGE);
        vm.loadProduct(1, choc, 5);
        vm.loadProduct(2, snack, 3);
        vm.loadProduct(3, bev,   2);
        vm.showInventory();

        // Scenario 1: Exact payment
        System.out.println("========== Scenario 1: Exact Payment ==========");
        vm.insertMoney(1.50);
        vm.selectProduct(1);

        // Scenario 2: Overpayment & Change
        System.out.println("========== Scenario 2: Overpayment & Change ==========");
        vm.insertMoney(3.00);
        vm.selectProduct(2);

        // Scenario 3: Underpayment & Refund
        System.out.println("========== Scenario 3: Underpayment & Refund ==========");
        vm.insertMoney(1.00);
        vm.selectProduct(3);

        // Scenario 4: Deplete Rack & Retry
        System.out.println("========== Scenario 4: Deplete Rack 3 & Retry ==========");
        vm.insertMoney(5.00); vm.selectProduct(3);
        vm.insertMoney(2.50); vm.selectProduct(3);
        vm.insertMoney(2.50); vm.selectProduct(3);

        vm.showInventory();
    }
}
\end{lstlisting}

\subsection{Wrapping up}\label{oduvT-U5GDbYO8Cu-X2TP}
\\
In this chapter, we've explored the complete design of a vending machine. We 've also examined how a basic vending machine can be visualized using various UML diagrams and designed using object-oriented principles and design patterns.

\subsubsection{}\label{DzQq279g7r9A1b-iCoiDe}

% End of section content